\documentclass[12pt,a4paper]{article}

\usepackage[utf8]{inputenc}
\usepackage[T1]{fontenc}
\usepackage[brazil]{babel}
\usepackage{amsmath,amssymb}
\usepackage{booktabs}
\usepackage{array}
\usepackage{geometry}

\geometry{
left=2cm,
right=2cm,
top=2cm,
bottom=2cm
}

\begin{document}

\begin{center}
{\LARGE \textbf{FOLHA DE REVISÃO - INTEGRAÇÃO NUMÉRICA}}
\end{center}

\hrule

\section*{1. Roteiro para Resolver Qualquer Questão}

Dada uma integral

\[
I=\int_a^b f(x)\,dx
\]

sempre siga esta sequência:

\begin{enumerate}
\item Identifique os limites \(a\) e \(b\)
\item Identifique o método pedido
\item Determine o número de subintervalos \(n\)
\item Calcule \(h\)
\item Encontre os pontos \(x_i\)
\item Calcule os valores \(f(x_i)\)
\item Monte a tabela
\item Substitua na fórmula
\end{enumerate}

\hrule

\section*{2. Cálculo de \(h\)}

A largura dos subintervalos é

\[
h=\frac{b-a}{n}
\]

onde

\[
a=\text{limite inferior}
\]

\[
b=\text{limite superior}
\]

\[
n=\text{número de subintervalos}
\]

\hrule

\section*{3. Como Encontrar os Pontos}

Após calcular \(h\):

\[
x_i=a+i\cdot h
\]

Exemplo:

\[
x_0=a
\]

\[
x_1=a+h
\]

\[
x_2=a+2h
\]

\[
x_3=a+3h
\]

\[
x_4=a+4h
\]

\hrule

\section*{4. Trapézio Simples}

Sempre utiliza:

\[
n=1
\]

Coeficientes:

\[
1 \qquad 1
\]

Fórmula:

\[
I
\approx
\frac{b-a}{2}
\left[f(a)+f(b)\right]
\]

\textbf{Observação:}

Utiliza apenas os extremos do intervalo.

\hrule

\section*{5. Trapézio Composto}

O valor de \(n\) é informado no exercício.

Coeficientes:

\[
1 \quad 2 \quad 2 \quad \cdots \quad 2 \quad 1
\]

Fórmula:

\[
I
\approx
\frac{h}{2}
\left[
f_0
+
2f_1
+
2f_2
+
\cdots
+
2f_{n-1}
+
f_n
\right]
\]

Exemplo para \(n=4\):

\[
1 \quad 2 \quad 2 \quad 2 \quad 1
\]

\hrule

\section*{6. Simpson 1/3}

Sempre utiliza:

\[
n=2
\]

Coeficientes:

\[
1 \qquad 4 \qquad 1
\]

Fórmula:

\[
I
\approx
\frac{h}{3}
\left[
f_0+4f_1+f_2
\right]
\]

Passos:

\[
h=\frac{b-a}{2}
\]

\[
x_0=a
\]

\[
x_1=a+h
\]

\[
x_2=b
\]

\textbf{Decore:}

\[
1 \qquad 4 \qquad 1
\]

\hrule

\section*{7. Simpson 3/8}

Sempre utiliza:

\[
n=3
\]

Coeficientes:

\[
1 \qquad 3 \qquad 3 \qquad 1
\]

Fórmula:

\[
I
\approx
\frac{3h}{8}
\left[
f_0+3f_1+3f_2+f_3
\right]
\]

Passos:

\[
h=\frac{b-a}{3}
\]

\[
x_0=a
\]

\[
x_1=a+h
\]

\[
x_2=a+2h
\]

\[
x_3=b
\]

\textbf{Decore:}

\[
1 \qquad 3 \qquad 3 \qquad 1
\]

\hrule

\section*{8. Exercício Resolvido}

Calcular

\[
I=
\int_0^{2\pi}
(6+3\cos x)\,dx
\]

Defina:

\[
f(x)=6+3\cos x
\]

\subsection*{Valores da Função}

\[
f(0)=9
\]

\[
f\left(\frac{\pi}{2}\right)=6
\]

\[
f(\pi)=3
\]

\[
f\left(\frac{3\pi}{2}\right)=6
\]

\[
f(2\pi)=9
\]

\begin{center}
\begin{tabular}{cc}
\toprule
$x$ & $f(x)$\\
\midrule
$0$ & $9$\\
$\frac{\pi}{2}$ & $6$\\
$\pi$ & $3$\\
$\frac{3\pi}{2}$ & $6$\\
$2\pi$ & $9$\\
\bottomrule
\end{tabular}
\end{center}

\hrule

\section*{9. Subintervalos da Questão}

\subsection*{Trapézio Simples}

\[
n=1
\]

\[
h=2\pi
\]

Pontos:

\[
0,\;2\pi
\]

\subsection*{Trapézio Composto (\(n=4\))}

\[
h=\frac{2\pi}{4}
=
\frac{\pi}{2}
\]

Pontos:

\[
0,\;
\frac{\pi}{2},\;
\pi,\;
\frac{3\pi}{2},\;
2\pi
\]

\subsection*{Simpson 1/3}

\[
n=2
\]

\[
h=\pi
\]

Pontos:

\[
0,\;
\pi,\;
2\pi
\]

\subsection*{Simpson 3/8}

\[
n=3
\]

\[
h=\frac{2\pi}{3}
\]

Pontos:

\[
0,\;
\frac{2\pi}{3},\;
\frac{4\pi}{3},\;
2\pi
\]

\hrule

\section*{10. Valores Trigonométricos Mais Importantes}

\[
\cos(0)=1
\]

\[
\cos\left(\frac{\pi}{2}\right)=0
\]

\[
\cos(\pi)=-1
\]

\[
\cos\left(\frac{3\pi}{2}\right)=0
\]

\[
\cos(2\pi)=1
\]

\hrule

\section*{11. Resumo para Decorar}

\begin{center}
\begin{tabular}{cc}
\toprule
Método & Coeficientes \\
\midrule
Trapézio Simples & $1 \;\; 1$ \\
Trapézio Composto & $1 \;\; 2 \;\; 2 \;\; \cdots \;\; 2 \;\; 1$ \\
Simpson 1/3 & $1 \;\; 4 \;\; 1$ \\
Simpson 3/8 & $1 \;\; 3 \;\; 3 \;\; 1$ \\
\bottomrule
\end{tabular}
\end{center}

\hrule

\section*{12. Erros Mais Comuns}

\begin{enumerate}
\item Esquecer de calcular \(h\)
\item Errar os pontos \(x_i\)
\item Utilizar coeficientes incorretos
\item Esquecer o fator da fórmula:
\[
\frac{h}{2},
\qquad
\frac{h}{3},
\qquad
\frac{3h}{8}
\]
\item Errar os valores trigonométricos
\end{enumerate}

\end{document}